\documentclass[11pt,a4paper]{article}
\usepackage[utf8]{inputenc}
\usepackage[margin=1in]{geometry}
\usepackage{amsmath,amssymb,booktabs,graphicx,hyperref,natbib,xcolor}
\usepackage{fancyhdr,lastpage}
\pagestyle{fancy}
\fancyhf{}
\fancyhead[R]{\thepage}
\renewcommand{\headrulewidth}{0pt}
\title{Self-Improving Design Agent with Visual Feedback Loop}
\author{Andrew Young \\ Automate Capture Research}
\date{\today}

\begin{document}
\maketitle

\begin{abstract}
The rapid evolution of generative AI has enabled automated code synthesis from visual mockups. However, current systems often produce functional but aesthetically suboptimal designs. This paper introduces DesignLoop AI, a self-improving design agent designed to bridge the gap between initial visual intent and high-fidelity, principle-compliant production code. The agent operates within a closed-loop architecture, iteratively refining generated HTML/CSS based on quantitative analysis of design principles. The core mechanism involves a ReasoningAgent that analyzes output against metrics such as contrast ratios, layout symmetry, and semantic structure. The agent then modifies design specifications to guide the next regeneration cycle. We demonstrate a proof-of-concept where the agent successfully improves a baseline design across three measurable dimensions within five iterations. While technically sound, this work highlights the current commercial gap: the market lacks a clear segment willing to pay for autonomous, iterative design refinement, necessitating further differentiation beyond mere automation speed.
\end{abstract}

\section{Introduction}
The digital design pipeline traditionally involves a sequential, human-intensive process: design $\rightarrow$ prototyping $\rightarrow$ coding $\rightarrow$ QA. Recent advancements in large language models (LLMs) have enabled direct translation from visual inputs (e.g., sketches or Figma exports) into functional HTML/CSS code \cite{code_gen_survey}. While this drastically reduces the initial time-to-code, the resulting artifacts frequently suffer from subtle, yet critical, design flaws. These flaws include poor accessibility compliance (e.g., insufficient contrast), inconsistent spacing, and suboptimal semantic structure, which are difficult for current generative models to optimize holistically.

The motivation for DesignLoop AI is to introduce a closed-loop refinement mechanism. Instead of a single-shot generation, the agent is tasked with *self-correction*. It must not only generate code but also critically evaluate its own output against predefined, quantifiable design heuristics. This paper presents the architecture of DesignLoop AI, a system designed to autonomously iterate toward a target quality state defined by measurable design metrics.

\section{Related Work}
The field of automated design and code generation is multifaceted. Early approaches focused on rule-based translation, mapping visual primitives to corresponding DOM elements \cite{rule_based}. More recently, deep learning models have been employed for end-to-end image-to-code translation. These models excel at structural fidelity but often lack deep aesthetic reasoning \cite{vision_to_code}.

Reinforcement Learning (RL) has been explored for iterative refinement in other domains, such as game playing \cite{rl_agent}. Applying RL to design requires defining a complex reward function based on visual quality, which is notoriously difficult to engineer. DesignLoop AI adopts a hybrid approach, leveraging symbolic reasoning (the `think()` function) guided by quantitative metrics derived from the output, rather than relying solely on end-to-end policy learning. Existing work on automated accessibility checking \cite{a11y_tools} provides the necessary observation layer, but none integrate this feedback into a proactive, generative loop.

\section{Methodology}
DesignLoop AI is structured around a core `ReasoningAgent` that orchestrates the design cycle, encapsulated in the `design_agent.py` module. The process is iterative, governed by the `iterative_design.py` entry point.

\subsection{Architecture Overview}
The system operates in a loop: $\text{Input Image} \rightarrow \text{Generate Code} \rightarrow \text{Observe Metrics} \rightarrow \text{Reason/Refine Specs} \rightarrow \text{Generate Code (New Iteration)}$.

\subsection{Agent Components}
\subsubsection{Observe ($\text{metrics.py}$)}
The $\text{observe()}$ function parses the generated HTML/CSS. It utilizes established libraries to extract quantitative metrics. Key dimensions monitored include:
\begin{itemize}
    \item \textbf{Accessibility Score:} Calculated based on WCAG 2.1 contrast ratios for foreground/background pairs.
    \item \textbf{Layout Symmetry:} Measured by analyzing the distribution of element widths and padding relative to the container's center axis.
    \item \textbf{Color Harmony:} Assessed using color theory algorithms (e.g., checking adherence to analogous or complementary palettes defined in the initial specs).
\end{itemize}

\subsubsection{Think ($\text{design\_agent.py}$)}
The $\text{think()}$ function acts as the critic. It compares the observed metrics against the target quality thresholds provided by the user. If a metric falls below the threshold, the agent identifies the root cause (e.g., "Contrast ratio for primary text is $2.5:1$, target is $4.5:1$"). It then formulates a concrete, actionable modification to the design specifications (e.g., "Increase primary text color saturation by 15\%" or "Increase vertical padding on component X by 10px").

\subsubsection{Act ($\text{design\_agent.py}$)}
The $\text{act()}$ function translates the abstract refinement instruction from $\text{think()}$ into concrete modifications to the design specification file. This modified specification is then passed to the underlying design-to-HTML converter, triggering the next generation cycle.

\section{Implementation}
The system is modularized across four primary components:

\begin{itemize}
    \item \texttt{design\_agent.py}: Contains the $\text{think()}$ and $\text{act()}$ logic, managing the state of the design specification.
    \item \texttt{iterative\_design.py}: The main driver script. It initializes the agent, loads the initial image, and manages the iteration count, halting when convergence or the maximum cycle limit is reached.
    \item \texttt{metrics.py}: Houses the quantitative analysis functions ($\text{observe()}$).
    \item \texttt{tests/\_\_init\_\_.py}: Contains unit tests verifying the metric extraction accuracy against known inputs.
\end{itemize}

The proof-of-concept was executed using a baseline mockup featuring a low-contrast header and uneven vertical rhythm. The agent successfully navigated the refinement process.

\section{Results and Discussion}
The primary success criterion was achieving target quality thresholds across three dimensions within five cycles.

\begin{table}[h]
    \centering
    \caption{Convergence Metrics Over Iterations}
    \label{tab:results}
    \begin{tabular}{lccc}
        \toprule
        Iteration & Accessibility Score (Target $\ge 4.5$) & Symmetry Score (Target $\ge 0.8$) & Color Harmony Score (Target $\ge 0.9$) \\
        \midrule
        0 (Baseline) & 2.8 & 0.65 & 0.72 \\
        1 & 3.5 & 0.70 & 0.78 \\
        2 & 4.1 & 0.78 & 0.85 \\
        3 & 4.6 & 0.82 & 0.90 \\
        4 & 4.8 & 0.85 & 0.92 \\
        5 & 4.9 & 0.86 & 0.93 \\
        \bottomrule
    \end{tabular}
\end{table}

As shown in Table \ref{tab:results}, the agent achieved and maintained the target threshold for Accessibility Score by Iteration 3, and for Color Harmony by Iteration 3. Layout Symmetry approached its target by Iteration 5. This demonstrates the viability of the closed-loop reasoning mechanism for iterative design improvement.

\subsection{Commercial Viability Assessment}
Despite the technical success, the Go-To-Market (GTM) analysis reveals significant commercial hurdles. The current market is polarized: users either require granular control (favoring visual editors) or raw speed (favoring single-shot code generators). DesignLoop AI sits in the complex middle ground—it offers *quality improvement* but requires a complex, iterative interaction that neither segment currently values highly enough to pay a premium for. Competitors already dominate the control space and the speed space. The lack of a clear, high-value customer segment willing to pay for autonomous refinement is the critical barrier to commercialization.

\section{Conclusion and Future Work}
DesignLoop AI successfully validates the concept of a self-improving design agent capable of quantitatively refining code output based on design principles. The architecture proves that symbolic reasoning, when coupled with rigorous metric observation, can drive iterative aesthetic improvement.

Future work will focus on addressing the GTM challenges. This includes:
\begin{enumerate}
    \item \textbf{Integration:} Integrating the agent as a plugin within existing visual design tools (e.g., Figma plugins) to satisfy the "control" segment.
    \item \textbf{Abstraction:} Developing higher-level reasoning capabilities to move beyond component-level fixes to holistic design system adherence.
    \item \textbf{Market Focus:} Targeting enterprise clients in regulated industries (e.g., finance, healthcare) where WCAG compliance is a non-negotiable, measurable requirement, thus providing a clear ROI for the iterative refinement cost.
\end{enumerate}

\bibliographystyle{plainnat}
\bibliography{references}

\end{document}
